% Section 6.2, from lectures/L06/README.md: the objectives, the evaluation questions, and what
% comes after the last lecture.

\section{Review}
\label{c6:sec:review}

\needspace{6\baselineskip}
\subsection*{What you should be able to do now}
After this chapter, you should be able to:
\begin{itemize}
  \item Explain what a thread is and how it is created using \code{std::thread}.
  \item Explain the difference between concurrency and parallelism.
  \item Explain preemptive and cooperative scheduling.
  \item Explain what a data race is and why it leads to \emph{undefined behavior}.
  \item Use \code{std::mutex} and \code{std::lock_guard} correctly.
  \item Explain when \code{std::atomic} is sufficient and when \code{std::mutex} is required.
  \item Use \code{std::condition_variable} to signal between threads without busy-waiting.
  \item Launch asynchronous tasks using \code{std::async} and retrieve results via
    \code{std::future}.
  \item Explain what priority inversion is and why \code{std::mutex} does not solve it on an RTOS.
\end{itemize}

\needspace{6\baselineskip}
\subsection*{Questions to test yourself}
You should be able to explain:
\begin{enumerate}
  \item What is the difference between a \code{mutex} and an \code{atomic}?
  \item What is a data race?
  \item What is the difference between concurrency and parallelism?
  \item When would you use a \code{condition_variable} instead of polling in a loop?
  \item What is priority inversion, and why should you be careful with \code{std::mutex} on an
    RTOS?
\end{enumerate}

\needspace{6\baselineskip}
\subsection*{Looking ahead}
This is the last chapter, and with it the book's main text is complete. What remains is to work
the exercises in \secref{c6:sec:exercises} and compare your programs with the solutions in the
companion repository. After that, the two written papers in
Appendices~\ref{exam:ch:about} to \ref{ansb:ch:answers} are there to test, on paper and without
a compiler, how much of all six chapters you can reconstruct on your own.
